\begin{table}[t]
  \centering
  \caption{\textbf{本稿の主張の境界。}
  実測した比較、実装済みの契約、
  これから検証する仮説、未測定の項目を分ける。
  }
  \label{tab:claim_scope}
  \small
  \begin{tabularx}{\linewidth}{@{}p{0.36\linewidth}p{0.14\linewidth}X@{}}
    \toprule
    主張 & 状態 & 根拠 \\
    \midrule
    失敗時に停止するメートル系アライメント & \statusimplemented &
      地上面推定、共通尺度の二面当てはめ、
      fit/holdout 分離、トポロジーと意味セグメントgate \\
    軌道生成とグループ分割 & \statusimplemented &
      凸包制約、弧長一様サンプル、trajectory group 分割 \\
    単一変換からの整合教師 & \statusimplemented &
      キーポイント・領域・線・線の意味・姿勢を同一変換から生成 \\
    検出器の並列姿勢回帰 & \statusimplemented &
      共有特徴からの密予測と姿勢10値の回帰 \\
    交差モデル比較（実写 validation / 合成 test） & \statusimplemented &
      domain ごと決定論的 120 件、共通 RANSAC 手続き、全推論 CPU \\
    視点拡張が実写汎化を改善する & \statushypothesis &
      同一枚数の狭域・広域比較と実写評価が必要 \\
    未知会場への汎化 & \statusunmeasured &
      会場を分けた評価集合が必要 \\
    3DGS 利用の因果効果 & \statusunmeasured &
      実写のみの同一アーキテクチャ比較が必要 \\
    採用視点分布の集計 & \statusunmeasured &
      学習へ渡った視点の方位・高さ・距離の集計図は未作成 \\
    \bottomrule
  \end{tabularx}
\end{table}
